%
% this file is encoded in utf-8
% v1.7
% 填入你的論文題目、姓名等資料
% 如果題目內有必須以數學模式表示的符號，請用 \mbox{} 包住數學模式，如下範例
% 如果中文名字是單名，與姓氏之間建議以全形空白填入，如下範例
% 英文名字中的稱謂，如 Prof. 以及 Dr.，其句點之後請以不斷行空白~代替一般空白，如下範例
% 如果你的指導教授沒有如預設的三位這麼多，則請把相對應的多餘教授的中文、英文名
%    的定義以空的大括號表示
%    如，\renewcommand\advisorCnameB{}
%          \renewcommand\advisorEnameB{}
%          \renewcommand\advisorCnameC{}
%          \renewcommand\advisorEnameC{}

% Use "~" instead of space bar for any space you want

% 論文題目 (中文)
% Chinese Thesis Title
\newcommand\cTitle{
    有效的語言習得和大型語言模型的認知 
}

% 論文題目 (英文)
% English Thesis Title
\newcommand\eTitle{
    Effective Language Acquisition and Cognition of Large Language Models
}

% 我的姓名 (中文)
% Chinese Name
\newcommand\myCname{林哲森}
% 我的姓名 (英文)
% English Name
\newcommand\myEname{Luis Frentzen Salim}
% 我的學號
% Student ID
\newcommand\myStudentID{M11315804}
% 指導教授A的姓名 (中文)
% Chinese Name for Advisor A
\newcommand\advisorCnameA{鮑興國~博士}
% 指導教授A的姓名 (英文)
% English Name for Advisor A
\newcommand\advisorEnameA{Dr. Hsing-Kuo~Pao}
% 指導教授B的姓名 (中文)
% Chinese Name for Advisor B
\newcommand\advisorCnameB{古倫維~博士}
% 指導教授B的姓名 (英文)
% English Name for Advisor B
\newcommand\advisorEnameB{Dr. Lun-Wei~Gu}
% 指導教授C — 無第三位，保持空白 (No third advisor)
\newcommand\advisorCnameC{}
\newcommand\advisorEnameC{}
% 校名 (中文)
\newcommand\univCname{國立台灣科技大學}
% 校名 (英文)
%\renewcommand\univEname{National Taiwan University of science and technology}

% 系所名 (中文)
\newcommand\deptCname{資訊工程系}

% 系所全名 (英文)
%\renewcommand\fulldeptEname{Graduate School of Electro-Optical Engineering}
% 系所短名 (英文, 用於書名頁學位名領域)
%\renewcommand\deptEname{Electro-Optical Engineering}
% 學院英文名 (如無，則以空的大括號表示)
%\renewcommand\collEname{College of Electrical and Communication Engineering}

% 學位名 (中文)
\newcommand\degreeCname{碩士}
% 學位名 (英文)
%\renewcommand\degreeEname{Master of Science}
% 口試年份 (中文、民國)
\newcommand\cYear{一一四}
% 口試月份 (中文)
\newcommand\cMonth{七} 
% 口試月份 (中文)
\newcommand\cDay{二十九} 
% 口試年份 (阿拉伯數字、西元)
%\renewcommand\eYear{2009} 
% 口試月份 (英文)
%\renewcommand\eMonth{July}
% 學校所在地 (英文)
%\renewcommand\ePlace{Taipei, Taiwan}
%畢業級別；用於書背列印；若無此需要可忽略
\newcommand\GraduationClass{95}

%%%%%%%%%%%%%%%%%%%%%%